% =====================================================================
% Fiche de revision - Python Scripting - Lesson 2 : POO
% Compilation : pdfLaTeX (Overleaf), aucune option particuliere requise.
% =====================================================================

\documentclass[a4paper,10pt]{article}

\usepackage[utf8]{inputenc}
\usepackage[T1]{fontenc}
\usepackage[french]{babel}
\usepackage[a4paper,margin=1.3cm]{geometry}
\usepackage{xcolor}
\usepackage{multicol}
\usepackage{booktabs}
\usepackage{array}
\usepackage{tabularx}
\usepackage{enumitem}
\usepackage{listings}
\usepackage{titlesec}
\usepackage{tcolorbox}

% ---------------------------------------------------------------- couleurs
\definecolor{accent}{HTML}{1F5FA8}
\definecolor{codefg}{HTML}{A03030}
\definecolor{codebg}{HTML}{F4F6F8}
\definecolor{kw}{HTML}{0B5FA5}
\definecolor{cmt}{HTML}{6C8A54}
\definecolor{strcol}{HTML}{2E7D32}
\definecolor{danger}{HTML}{B3261E}

% ---------------------------------------------------------------- titres
\titleformat{\section}
  {\normalfont\large\bfseries\color{accent}}{}{0pt}{}
  [\vspace{-0.75em}{\color{accent}\rule{\linewidth}{0.7pt}}]
\titlespacing*{\section}{0pt}{1.4ex}{0.7ex}

\titleformat{\subsection}[runin]
  {\normalfont\bfseries}{}{0pt}{}[\;---\;]
\titlespacing*{\subsection}{0pt}{0.8ex}{0.4em}

% ---------------------------------------------------------------- code
\lstset{
  language=Python,
  basicstyle=\ttfamily\footnotesize,
  keywordstyle=\color{kw}\bfseries,
  commentstyle=\color{cmt}\itshape,
  stringstyle=\color{strcol},
  showstringspaces=false,
  columns=fullflexible,
  keepspaces=true,
  backgroundcolor=\color{codebg},
  frame=none,
  xleftmargin=0.3em,
  aboveskip=0.55em,
  belowskip=0.55em,
  breaklines=true,
}

% code en ligne
\newcommand{\co}[1]{\texttt{\textcolor{codefg}{#1}}}

% ---------------------------------------------------------------- boites
\newtcolorbox{piegebox}[1]{
  colback=red!3, colframe=danger, boxrule=0.7pt,
  left=5pt, right=5pt, top=4pt, bottom=4pt,
  title=#1, fonttitle=\bfseries\small, coltitle=white
}

\newtcolorbox{astucebox}[1]{
  colback=accent!5, colframe=accent, boxrule=0.7pt,
  left=5pt, right=5pt, top=4pt, bottom=4pt,
  title=#1, fonttitle=\bfseries\small, coltitle=white
}

\setlist[itemize]{leftmargin=1.1em, itemsep=0.15em, topsep=0.25em, parsep=0pt}
\setlength{\columnsep}{1.1em}
\setlength{\parindent}{0pt}
\renewcommand{\arraystretch}{1.15}

\pagestyle{empty}

\begin{document}

% Desactive les espaces automatiques du francais avant : ; ! ?
% pour que le code en ligne (ex. "def f(self):") ne soit pas deforme.
\shorthandoff{;:!?}

\small

\begin{center}
  {\LARGE\bfseries\color{accent} Python Scripting --- Lesson 2}\\[0.3em]
  {\normalsize Programmation orientée objet \quad$\cdot$\quad fiche de révision QCM}
\end{center}
\vspace{0.4em}

% =====================================================================
\begin{piegebox}{Les 16 réponses pièges à connaître par c\oe ur}
\begin{minipage}[t]{0.485\linewidth}
\begin{tabularx}{\linewidth}{@{}l>{\raggedright\arraybackslash}X@{}}
\toprule
\textbf{Situation} & \textbf{Résultat} \\
\midrule
accès par défaut          & \textbf{public} --- aucun mot-clé en Python \\
\co{\_x}                  & \textbf{convention} seule, rien n'est bloqué \\
\co{t.\_\_x} (privé)      & \co{AttributeError} \\
\co{t.\_C\_\_x}           & marche --- \emph{name mangling} \\
\co{\_\_init\_\_} qui \co{return 5} & \co{TypeError} --- doit rendre \co{None} \\
enfant sans \co{super()}  & attributs du parent \textbf{absents} \\
\co{type(d) is Animal}    & \co{False} pour \co{d = Dog()} \\
\co{isinstance(d, Animal)}& \co{True} --- suit l'héritage \\
\bottomrule
\end{tabularx}
\end{minipage}\hfill
\begin{minipage}[t]{0.485\linewidth}
\begin{tabularx}{\linewidth}{@{}l>{\raggedright\arraybackslash}X@{}}
\toprule
\textbf{Situation} & \textbf{Résultat} \\
\midrule
\co{a + b} sans \co{\_\_add\_\_}  & \co{TypeError} \\
\co{3 * v} sans \co{\_\_rmul\_\_} & \co{TypeError} (\co{v * 3} passe) \\
\co{+=} sans \co{\_\_iadd\_\_}    & retombe sur \co{\_\_add\_\_} \\
\co{!=} sans \co{\_\_ne\_\_}      & déduit de \co{\_\_eq\_\_} (Python 3) \\
\co{\_\_eq\_\_} sans \co{\_\_hash\_\_} & objet \textbf{non hachable} \\
\co{print(obj)} sans \co{\_\_str\_\_} & utilise \co{\_\_repr\_\_} \\
ni l'un ni l'autre        & \co{<\_\_main\_\_.P object at 0x...>} \\
\co{\_\_div\_\_}          & n'existe pas : c'est \co{\_\_truediv\_\_} \\
\bottomrule
\end{tabularx}
\end{minipage}
\end{piegebox}

\vspace{0.5em}
\begin{multicols}{2}

% =====================================================================
\section{Les concepts et leur nom exact}

\begin{tabularx}{\linewidth}{@{}l>{\raggedright\arraybackslash}X@{}}
\toprule
\textbf{Class}         & le modèle : attributs (données) et méthodes (comportement) \\
\textbf{Object}        & une \emph{instance} concrète de la classe \\
\textbf{Encapsulation} & regrouper données \textbf{et} méthodes dans une seule entité \\
\textbf{Abstraction}   & cacher les détails, n'exposer que l'essentiel \\
\textbf{Inheritance}   & une classe reprend attributs et méthodes d'une autre \\
\textbf{Polymorphism}  & même nom de méthode, comportement différent selon la classe \\
\textbf{Composition}   & un attribut est lui-même l'objet d'une autre classe \\
\bottomrule
\end{tabularx}

\vspace{0.4em}
\textbf{Le piège classique du QCM} : une question donne six définitions et
mélange les sept mots. Encapsulation $\neq$ abstraction : la première
\emph{regroupe}, la seconde \emph{cache}.

% =====================================================================
\section{Classe, objet, instanciation}

\begin{lstlisting}
class Person:
    def __init__(self, name, age):
        self.name = name
        self.age = age

    def greet(self):
        print(f"Hello, {self.name}")

p = Person("Alice", 30)   # pas de "new"
p.greet()
\end{lstlisting}

Appeler la classe suffit : Python crée l'objet puis exécute
\co{\_\_init\_\_}. Il n'y a \textbf{pas} de mot-clé \co{new}.

\subsection{Connaître la classe d'un objet}
\co{type(p)} donne la classe, \co{type(p).\_\_name\_\_} son nom en
\textbf{chaîne} (\co{'Person'}) --- pratique dans un \co{\_\_repr\_\_}
générique.

\subsection{Attribut de classe contre attribut d'instance}
\begin{lstlisting}
class Counter:
    total = 0          # de CLASSE : partage
    def __init__(self):
        self.count = 0 # d'INSTANCE : propre

Counter.total = 5      # vu par a.total ET b.total
\end{lstlisting}
Un attribut défini dans \co{\_\_init\_\_} n'appartient qu'à cette instance ;
un attribut de classe est \textbf{partagé} par toutes. Conséquence :
\co{Dog.name} lève une \co{AttributeError} si \co{name} est posé dans
\co{\_\_init\_\_} --- il vit sur l'objet, pas sur la classe.

% =====================================================================
\section{\co{\_\_init\_\_} et \co{self}}

\begin{itemize}
  \item \co{\_\_init\_\_} est appelé \textbf{automatiquement} à chaque
        instanciation ; il \emph{initialise} l'objet, il ne le crée pas.
  \item Il doit renvoyer \co{None} : tout autre \co{return} lève une
        \co{TypeError}.
  \item \co{self} n'est \textbf{pas} un mot-clé, c'est une convention :
        seule sa \textbf{position} (premier paramètre) compte.
  \item On ne passe \textbf{jamais} \co{self} à la main.
\end{itemize}

\begin{lstlisting}
p.greet()          # les deux sont
Person.greet(p)    # equivalents
\end{lstlisting}

Un paramètre avec valeur par défaut rend l'argument optionnel :
\co{def \_\_init\_\_(self, start=0)} accepte \co{Counter()} comme
\co{Counter(10)}.

% =====================================================================
\section{Encapsulation}

\begin{lstlisting}
class Test:
    varPublic     = 10  # public
    _varProtected = 20  # protege : 1 underscore
    __varPrivate  = 30  # prive   : 2 underscores
\end{lstlisting}

\begin{tabularx}{\linewidth}{@{}l>{\raggedright\arraybackslash}X@{}}
\toprule
\co{x}      & public --- le \textbf{défaut} en Python \\
\co{\_x}    & protégé : \textbf{simple convention}, rien n'est empêché \\
\co{\_\_x}  & privé : Python \textbf{renomme} en \co{\_Classe\_\_x} \\
\bottomrule
\end{tabularx}

\subsection{Name mangling}
\co{\_\_varPrivate} devient \co{\_Test\_\_varPrivate}. Donc
\co{t.\_\_varPrivate} lève une \co{AttributeError}, mais
\co{t.\_Test\_\_varPrivate} renvoie \co{30} : la valeur est
\textbf{cachée, pas protégée}. Son vrai but est d'éviter les collisions de
noms entre une classe et ses filles.

% =====================================================================
\section{Héritage}

\begin{lstlisting}
class Dog(Animal):          # Dog = sous-classe
    def __init__(self, name, breed):
        super().__init__(name)  # INDISPENSABLE
        self.breed = breed

    def speak(self):        # method overriding
        return "Woof!"
\end{lstlisting}

\begin{itemize}
  \item \co{Dog} : \emph{subclass} / \emph{child} ; \co{Animal} :
        \emph{superclass} / \emph{parent}.
  \item L'enfant peut \textbf{ajouter} attributs et méthodes, et
        \textbf{redéfinir} celles du parent (\emph{overriding}).
  \item Définir \co{\_\_init\_\_} dans l'enfant \textbf{remplace} celui du
        parent : sans \co{super().\_\_init\_\_()}, les attributs du parent
        n'existent pas et l'\co{AttributeError} tombe plus tard.
\end{itemize}

\subsection{Appeler la version du parent}
\co{super().speak()} ou \co{Animal.speak(self)}. \textbf{Jamais}
\co{self.speak()} : cela rappellerait la version de l'enfant --- récursion
infinie.

\subsection{Tester l'appartenance}
\co{isinstance(d, Animal)} vaut \co{True} (l'héritage compte) et
\co{issubclass(Dog, Animal)} aussi, mais \co{type(d)} vaut
\textbf{exactement} \co{Dog}.

% =====================================================================
\section{Polymorphisme}

\begin{lstlisting}
class Dog(Animal):
    def speak(self):
        return "Woof!"

class Bird(Animal):
    def speak(self, repeat=2):
        return "Tweet!" * repeat

for a in [Dog("Rex"), Bird("Tweety")]:
    print(a.speak())
# Woof!
# Tweet!Tweet!
\end{lstlisting}

Python choisit la méthode \textbf{à l'exécution}, d'après la classe réelle
de l'objet : la recherche part de la classe de l'objet puis remonte la
chaîne d'héritage. Le code appelant n'a pas besoin de connaître le type
exact.

\subsection{Autres exemples valables}
\co{len()} sur une chaîne, une liste ou un dictionnaire : un seul nom, une
implémentation par type (\co{\_\_len\_\_}). De même
\co{sum(s.computeArea() for s in shapes)} sur une liste mélangeant
rectangles, cercles et carrés.

% =====================================================================
\section{Composition et abstraction}

\begin{tabularx}{\linewidth}{@{}l>{\raggedright\arraybackslash}X@{}}
\toprule
\textbf{Héritage}    & relation \textbf{\emph{is-a}} : un \co{Square} \emph{est un} \co{Rectangle} \\
\textbf{Composition} & relation \textbf{\emph{has-a}} : un compte \emph{a} des transactions \\
\bottomrule
\end{tabularx}

\vspace{0.4em}
\textbf{Abstraction} = cacher l'implémentation et n'exposer que
l'essentiel : on appelle \co{deposit()} sans savoir comment l'historique est
stocké.

\subsection{Forcer une fille à implémenter une méthode}
\begin{lstlisting}
# 1. echoue a l'APPEL
def withdrawal(self):
    raise NotImplementedError

# 2. echoue a l'INSTANCIATION
from abc import ABC, abstractmethod

class Account(ABC):
    @abstractmethod
    def withdrawal(self): ...
\end{lstlisting}

Dans le cours, le parent écrit \co{def speak(self): pass} : c'est un
comportement par défaut à redéfinir, mais \co{pass} renvoie \co{None}
silencieusement, donc un oubli passe inaperçu.

\subsection{UML}
Un nom de classe ou de méthode \textbf{en italique} signifie
\textbf{abstrait} : la classe n'est pas faite pour être instanciée
directement.

% =====================================================================
\section{Surcharge des opérateurs}

\begin{lstlisting}
class Vector:
    def __init__(self, x, y, z):
        self.x, self.y, self.z = x, y, z

    def __add__(self, other):
        return Vector(self.x + other.x,
                      self.y + other.y,
                      self.z + other.z)

    def __mul__(self, scalar):
        return Vector(self.x * scalar,
                      self.y * scalar,
                      self.z * scalar)

    def __repr__(self):
        return f'Vector({self.x}, {self.y}, {self.z})'

print(Vector(1,2,3) + Vector(4,5,6))  # Vector(5, 7, 9)
print(Vector(1,2,3) * 3)              # Vector(3, 6, 9)
\end{lstlisting}

\begin{minipage}[t]{0.48\linewidth}
\begin{tabularx}{\linewidth}{@{}lX@{}}
\toprule
\multicolumn{2}{@{}l}{\textbf{Arithmétiques}} \\
\midrule
\co{+}  & \co{\_\_add\_\_} \\
\co{-}  & \co{\_\_sub\_\_} \\
\co{*}  & \co{\_\_mul\_\_} \\
\co{/}  & \co{\_\_truediv\_\_} \\
\co{//} & \co{\_\_floordiv\_\_} \\
\co{\%} & \co{\_\_mod\_\_} \\
\co{**} & \co{\_\_pow\_\_} \\
\bottomrule
\end{tabularx}
\end{minipage}\hfill
\begin{minipage}[t]{0.48\linewidth}
\begin{tabularx}{\linewidth}{@{}lX@{}}
\toprule
\multicolumn{2}{@{}l}{\textbf{Comparaison}} \\
\midrule
\co{<}  & \co{\_\_lt\_\_} \\
\co{>}  & \co{\_\_gt\_\_} \\
\co{<=} & \co{\_\_le\_\_} \\
\co{>=} & \co{\_\_ge\_\_} \\
\co{==} & \co{\_\_eq\_\_} \\
\co{!=} & \co{\_\_ne\_\_} \\
\bottomrule
\end{tabularx}
\end{minipage}

\vspace{0.5em}
\textbf{Affectations augmentées} : le \co{i} veut dire \emph{in place}.
\co{-=} \co{\_\_isub\_\_}, \co{*=} \co{\_\_imul\_\_}, \co{/=}
\co{\_\_itruediv\_\_}, \co{//=} \co{\_\_ifloordiv\_\_}, \co{\%=}
\co{\_\_imod\_\_}, \co{**=} \co{\_\_ipow\_\_}.

\subsection{Les quatre règles qui tombent}
\begin{itemize}
  \item Sans \co{\_\_add\_\_}, \co{obj1 + obj2} lève
        \co{TypeError: unsupported operand type(s)}.
  \item \co{\_\_add\_\_} doit renvoyer un \textbf{nouvel objet de la même
        classe} pour enchaîner \co{v1 + v2 + v3}.
  \item \co{v * 3} utilise \co{\_\_mul\_\_} ; \co{3 * v} exige la version
        \emph{réfléchie} \co{\_\_rmul\_\_} (l'objet est à droite).
  \item Sans \co{\_\_iadd\_\_}, \co{obj += other} devient
        \co{obj = obj + other} : le nom pointe sur un objet \textbf{neuf}.
        Définir \co{\_\_iadd\_\_} modifie \textbf{sur place} --- c'est ce
        qui rend \co{list += } économique.
\end{itemize}

\subsection{Comparaisons : les cas dérivés}
\begin{itemize}
  \item \co{\_\_eq\_\_} défini seul : \co{!=} en est déduit automatiquement.
  \item \co{\_\_eq\_\_} sans \co{\_\_hash\_\_} : les instances deviennent
        \textbf{non hachables}, donc inutilisables comme clés de
        dictionnaire ou éléments d'un \co{set}.
  \item \co{\_\_lt\_\_} défini seul : \co{a < b} marche, et \co{a > b}
        aussi --- Python tente \co{a.\_\_gt\_\_(b)}, échoue, puis bascule
        sur l'appel réfléchi \co{b.\_\_lt\_\_(a)}.
\end{itemize}

% =====================================================================
\section{\co{\_\_repr\_\_} et \co{\_\_str\_\_}}

\begin{itemize}
  \item \co{print(obj)} utilise \co{\_\_str\_\_} et \textbf{retombe sur}
        \co{\_\_repr\_\_} s'il n'existe pas. Définir le seul
        \co{\_\_repr\_\_} suffit donc pour \co{print()}, \co{str()} et
        \co{repr()}.
  \item Sans aucun des deux :
        \co{<\_\_main\_\_.Point object at 0x000...>} --- nom de classe et
        adresse mémoire.
  \item Un bon \co{\_\_repr\_\_} est \textbf{non ambigu} et permettrait
        idéalement de \textbf{recréer} l'objet : \co{Vector(5, 7, 9)}.
\end{itemize}

% =====================================================================
\section{Réflexes des exercices}

\begin{itemize}
  \item \textbf{Banque} : \co{BankAccount} parent abstrait,
        \co{CheckingAccount} et \co{SavingsAccount} redéfinissent
        \co{withdrawal()} --- \emph{overriding} et polymorphisme.
  \item \textbf{Historique} : le compte contient une \textbf{liste d'objets}
        \co{Transaction} (composition) et ne reste fiable que parce que
        seuls \co{deposit()} et \co{withdrawal()} modifient \co{\_balance}
        (encapsulation).
  \item \textbf{Accesseur} : \co{\_balance} plus \co{get\_balance()} force
        tout changement à passer par les méthodes, donc à être journalisé.
  \item \textbf{Formes} : \co{class Square(Rectangle)} plutôt que
        \co{Shape} --- un carré \emph{est un} rectangle, il réutilise
        largeur, hauteur et \co{computeArea()} en passant la même valeur
        deux fois à \co{super().\_\_init\_\_()}.
  \item \textbf{Fraction} : réduire par le \co{gcd} dans \co{\_\_init\_\_},
        sinon \co{Fractional(2, 4)} et \co{Fractional(1, 2)} ne sont pas
        égales pour un \co{\_\_eq\_\_} naïf.
  \item \textbf{Dénominateur nul} : \co{Fractional(42, 0)} doit
        \textbf{lever une exception} dès la construction, là où l'erreur est
        commise.
  \item \textbf{Comparer deux fractions} : \co{a*d < c*b} plutôt que
        \co{a/b < c/d} --- on reste en entiers, sans arrondi flottant
        (valable tant que les dénominateurs sont positifs, d'où la
        normalisation du signe).
\end{itemize}

\end{multicols}

\vspace{0.3em}
\begin{astucebox}{Checklist des 60 dernières secondes}
\begin{minipage}[t]{0.485\linewidth}
\begin{itemize}
  \item Tout est \textbf{public} par défaut ; \co{\_} = convention,
        \co{\_\_} = renommage.
  \item \co{\_\_init\_\_} ne renvoie \textbf{rien}, et \co{self} est une
        convention de position.
  \item Enfant avec \co{\_\_init\_\_} : \textbf{toujours}
        \co{super().\_\_init\_\_(...)}.
  \item \emph{is-a} $\to$ héritage ; \emph{has-a} $\to$ composition.
\end{itemize}
\end{minipage}\hfill
\begin{minipage}[t]{0.485\linewidth}
\begin{itemize}
  \item \co{/} donne \co{\_\_truediv\_\_} (jamais \co{\_\_div\_\_}) ;
        \co{i} veut dire \emph{in place}.
  \item \co{\_\_add\_\_} renvoie un \textbf{nouvel} objet de la classe.
  \item \co{\_\_eq\_\_} seul : \co{!=} suit, mais le \co{hash} disparaît.
  \item \co{\_\_repr\_\_} seul suffit pour \co{print}, \co{str} et
        \co{repr}.
\end{itemize}
\end{minipage}
\end{astucebox}

\vspace{0.4em}
\begin{astucebox}{Vocabulaire des énoncés (le QCM est en anglais)}
\begin{minipage}[t]{0.485\linewidth}
\begin{itemize}
  \item \emph{template} : modèle --- \emph{instance} : objet créé
  \item \emph{to instantiate} : instancier
  \item \emph{subclass / child} face à \emph{superclass / parent}
  \item \emph{to override} : redéfinir (et non « surcharger »)
  \item \emph{to overload an operator} : surcharger un opérateur
\end{itemize}
\end{minipage}\hfill
\begin{minipage}[t]{0.485\linewidth}
\begin{itemize}
  \item \emph{magic} ou \emph{dunder method} : méthode \co{\_\_x\_\_}
  \item \emph{name mangling} : renommage des membres privés
  \item \emph{unhashable} : inutilisable comme clé
  \item \emph{in place} : sur place, sans créer d'objet
  \item \emph{reflected} : version droite (\co{\_\_radd\_\_})
\end{itemize}
\end{minipage}
\end{astucebox}

\vspace{0.4em}
\begin{center}
\footnotesize\itshape
Format du contrôle : 6 propositions, 1 ou 2 bonnes réponses, une case fausse
annule la question.\\
Pour s'entraîner : \texttt{cd "Lesson 2" \&\& python quiz.py}
\end{center}

\end{document}
